\documentclass[11pt]{article}

\usepackage[margin=1in]{geometry}
\usepackage{amsmath,amssymb,amsthm}
\usepackage{hyperref}
\usepackage{enumitem}
\usepackage[numbers,sort&compress]{natbib}
\usepackage{hyperref}

\hypersetup{
  colorlinks=true,
  linkcolor=blue,
  citecolor=blue,
  urlcolor=blue
}

\newtheorem{lemma}{Lemma}
\newtheorem{exercise}{Exercise}

\title{Eternal Bang Cosmology \\ as a Pedagogical Case Study}
\author{A research note for: \\ James R. Van Dyke}
\date{\today}

\begin{document}
\maketitle

\begin{abstract}
We construct a multi-universe toy framework assembled from a finite collection of independent Friedmann--Lema\^{\i}tre--Robertson--Walker (FLRW) solutions and define global ``ensemble statistics'' such as a geometric-mean scale factor and an averaged Hubble rate. This bookkeeping permits simultaneous expansion and collapse among ensemble members and can be \emph{closed} by \emph{postulating} an FLRW-like evolution equation for the statistics, producing an ``Eternal Bang'' narrative of global cycles. We emphasize that this narrative does not emerge dynamically from general relativity unless two non-negotiable requirements are met: (i) a meta-structure that gives physical meaning to cross-universe time identification and any notion of ``separation,'' and (ii) covariant (Bianchi-consistent) completion or derivation of any inter-sector coupling introduced into a single universe's Friedmann equation.
\end{abstract}

\noindent\textbf{Keywords:} FLRW cosmology; toy models; averaging; Bianchi identity; effective fluids; pedagogy

\section{Introduction: toy models as conceptual laboratories}

Toy models are indispensable in cosmology: they build intuition, isolate mechanisms, and expose hidden assumptions \citep{Weinberg2008,Mukhanov2005,Dodelson2003}.
However, toy models also invite a common methodological failure: consistent bookkeeping can be mistaken for physical derivation, and imposed closure conditions can be mistaken for emergent dynamics.

This paper presents an explicit case study of that boundary. We assemble an ensemble of FLRW solutions and define global \emph{statistics} (e.g.\ a geometric-mean scale factor) that enable narratives such as simultaneous expansion and collapse and even global cycles. The central lesson is not that such constructions are invalid, but that they fall into distinct genres:
(i) a collection of independent solutions (kinematics),
(ii) a closed effective model (phenomenology),
or (iii) an embedded physical theory (dynamics).
Which genre applies depends on what is \emph{defined} versus what is \emph{derived}.

\paragraph{Context: averaging and the ``fitting problem.''}
The distinction between local dynamics and averaged/global descriptions is not unique to multi-universe parables; it is already central in the averaging/backreaction and fitting problems in cosmology \citep{Buchert2000,ClarksonEtAl2011}. Those works emphasize that averaging procedures and effective evolution laws require explicit choices (domains, measures, foliations) and that such choices do not automatically inherit the form of FLRW dynamics.

\medskip
\noindent\fbox{\parbox{\linewidth}{
\textbf{Checklist (parable $\rightarrow$ physics).} To elevate ensemble bookkeeping to a physical model one must (1) specify a \emph{meta-structure} that makes any cross-universe time identification meaningful, (2) ensure any modifications of the Friedmann constraint are \emph{covariantly completed} (Bianchi-consistent) or derived from an action, and (3) produce at least one falsifiable prediction not trivially mimicked by a simpler model.
}}

\section{Mathematical construction: ensemble and statistics}

\subsection{Ensemble of FLRW universes}

Let
\[
\mathcal{U}=\{U_i \mid i=0,1,\dots,N-1\},
\]
where each $U_i$ is modeled by an FLRW metric with its own cosmic time $t_i$:
\begin{equation}
\mathrm{d}s_i^2=-c^2\,\mathrm{d}t_i^2+a_i(t_i)^2\left[\frac{\mathrm{d}r^2}{1-k_i r^2}+r^2(\mathrm{d}\theta^2+\sin^2\theta\,\mathrm{d}\phi^2)\right],
\end{equation}
with $k_i\in\{-1,0,+1\}$.
Each universe contains a perfect fluid plus (optionally) a cosmological constant $\Lambda_i$, yielding the standard Friedmann equations \citep{Weinberg2008,Mukhanov2005,Dodelson2003}:
\begin{align}
H_i(t_i)^2 &\equiv \left(\frac{1}{a_i}\frac{\mathrm{d}a_i}{\mathrm{d}t_i}\right)^2
= \frac{8\pi G}{3}\rho_i(t_i) - \frac{k_i c^2}{a_i(t_i)^2} + \frac{\Lambda_i}{3}, \label{eq:localFriedmann}\\[4pt]
\frac{1}{a_i}\frac{\mathrm{d}^2 a_i}{\mathrm{d}t_i^2}
&= -\frac{4\pi G}{3}\!\left(\rho_i(t_i)+\frac{3p_i(t_i)}{c^2}\right)+\frac{\Lambda_i}{3}.
\end{align}
Local stress-energy conservation gives
\begin{equation}
\frac{\mathrm{d}\rho_i}{\mathrm{d}t_i} + 3H_i\!\left(\rho_i+\frac{p_i}{c^2}\right)=0. \label{eq:localConservation}
\end{equation}
The linkage among \eqref{eq:localFriedmann} and \eqref{eq:localConservation} is ultimately enforced by the contracted Bianchi identity in each spacetime \citep{Wald1984}.

\subsection{A chosen time identification and dot convention}

To discuss ``simultaneity across universes'' we must \emph{choose} an identification map. Concretely, fix functions $t_i(t)$ (monotone in $t$) and consider the induced ensemble of quantities along the parameter $t$:
\[
a_i(t) := a_i(t_i(t)),\qquad \rho_i(t):=\rho_i(t_i(t)),\qquad \text{etc.}
\]
\textbf{Convention:} from here on, a dot denotes $\mathrm{d}/\mathrm{d}t$ after a choice of all maps $t\mapsto t_i(t)$ has been made. This is a modeling decision, not a consequence of general relativity.

\subsection{Ensemble statistics}

Define the \emph{geometric-mean scale factor}
\begin{equation}
A(t):=\left(\prod_{i=0}^{N-1} a_i(t)\right)^{1/N},
\end{equation}
and the associated \emph{averaged Hubble rate}
\begin{equation}
H_A(t):=\frac{\dot A}{A}=\frac{1}{N}\sum_{i=0}^{N-1}\frac{\dot a_i}{a_i}.
\end{equation}
Also define an averaged density with weights $w_i\ge0$ and $\sum_i w_i=1$:
\begin{equation}
\rho_{\mathrm{eff}}(t):=\sum_{i=0}^{N-1} w_i\,\rho_i(t).
\end{equation}
These are \emph{statistics}. Absent additional structure, $A(t)$ is not the scale factor of any actual spacetime, and different choices of weights $w_i$ correspond to different, generally inequivalent ``measures'' on the ensemble \citep{Buchert2000,ClarksonEtAl2011}.

\section{The ``Eternal Bang'' narrative: sign mixing and imposed closure}

\subsection{Simultaneous expansion and collapse}

At a given parameter value $t$, define expansion or contraction of member $U_i$ by the sign of
\[
H_i(t):=\frac{\dot a_i(t)}{a_i(t)}.
\]
The ensemble exhibits \emph{simultaneous expansion and collapse} at $t$ if there exist $i,j$ such that $H_i(t)>0$ and $H_j(t)<0$.

A transparent construction uses mirror pairs. Assume $N$ is even and pair indices $i\leftrightarrow \tilde i$ so that
\begin{equation}
a_{\tilde i}(t)=a_i(-t),\qquad \rho_{\tilde i}(t)=\rho_i(-t),\qquad p_{\tilde i}(t)=p_i(-t).
\end{equation}
Then $H_{\tilde i}(t)=-H_i(-t)$, guaranteeing sign-mixing for $t\neq 0$ by construction.

\subsection{Global cycles via a postulated closure}

The statistics $(A,H_A,\rho_{\mathrm{eff}})$ do not satisfy any evolution law \emph{unless one is specified}. A common closure is to \emph{postulate} an FLRW-like constraint for $A(t)$:
\begin{equation}
H_A(t)^2 = \frac{8\pi G}{3}\rho_{\mathrm{eff}}(t) - \frac{Kc^2}{A(t)^2} + \frac{\Lambda_{\mathrm{eff}}}{3},
\label{eq:globalClosure}
\end{equation}
where $K\in\{-1,0,+1\}$ and $\Lambda_{\mathrm{eff}}$ are constants chosen at the ensemble level.
Selecting $K=+1$ and (for illustration) $\Lambda_{\mathrm{eff}}=0$ yields the standard closed-dust parametric form for $A(t)$ \citep{Dodelson2003,Mukhanov2005}.

\paragraph{Important limitation.}
Any cyclic ``Eternal Bang'' narrative obtained by repeating or gluing solutions of \eqref{eq:globalClosure} is purely a feature of the \emph{imposed} closure. We do not address singularity resolution at $A\to0$ or entropy accumulation across cycles, which require physics beyond this kinematic construction.

\section{Two critical gates}

\subsection{Gate 1: the meta-structure requirement}

In general relativity, cosmic time is defined within a spacetime (e.g.\ proper time of comoving observers) and is foliation-dependent in general \citep{Wald1984}. For disconnected manifolds there is no canonical identification $t_i\leftrightarrow t_j$.
Thus statements like ``at time $t$ some universes expand while others contract'' are physically meaningful only \emph{relative to} a chosen meta-structure that selects (or justifies) a cross-universe identification.

Examples of meta-structures include:
\begin{itemize}[leftmargin=1.2em]
\item A higher-dimensional embedding (e.g.\ branes in a bulk) where a bulk time defines the identification.
\item Minisuperspace quantum cosmology, where $\{a_i\}$ are degrees of freedom of a single system and time emerges relationally.
\item A reformulation in which the $U_i$ label widely separated regions in one manifold, connecting the discussion to averaging/fitting issues \citep{Buchert2000,ClarksonEtAl2011}.
\end{itemize}
Without such structure, $t$ and interpretations of $A(t)$ as ``separation'' remain conventions.

\subsection{Gate 2: the Bianchi requirement (covariant completion)}

Within a single FLRW spacetime, the Friedmann constraint, the acceleration equation, and the continuity equation are linked by the contracted Bianchi identity \citep{Wald1984}. Consequently, modifying a Friedmann constraint by adding an externally prescribed term does not define a consistent model unless one of the following is done:
\begin{enumerate}[leftmargin=1.4em]
\item \textbf{Effective-fluid completion:} reinterpret the added term as a conserved stress-energy component with density $\rho_X$ and pressure $p_X$ satisfying $\nabla_\mu T_X^{\mu\nu}=0$ in that spacetime; or
\item \textbf{Parent-theory derivation:} derive the modification from a covariant action (e.g.\ extra fields, modified gravity, embedding), so that Bianchi-consistency is automatic.
\end{enumerate}

Under effective-fluid completion in an FLRW background, conservation implies
\begin{equation}
\dot\rho_X + 3H(\rho_X+p_X/c^2)=0
\quad\Longrightarrow\quad
w_X(a):=\frac{p_X}{\rho_X c^2}
= -1-\frac{1}{3}\frac{\mathrm{d}\ln\rho_X}{\mathrm{d}\ln a},
\label{eq:wX}
\end{equation}
so any externally specified $\rho_X(a)$ corresponds observationally to a dark component with an implied equation of state.

\subsection{Lemma: averaging $\neq$ scaling}

\begin{lemma}[Averaging does not preserve dust scaling]
Let $a_i(t)>0$ for $i=0,\dots,N-1$ and define
\[
A(t)=\left(\prod_{i=0}^{N-1} a_i(t)\right)^{1/N}.
\]
Then for any fixed $t$,
\[
\frac{1}{N}\sum_{i=0}^{N-1} a_i(t)^{-3}\;\ge\; A(t)^{-3},
\]
with equality if and only if $a_0(t)=a_1(t)=\cdots=a_{N-1}(t)$.

Consequently, if each universe is dust-dominated with $\rho_i(t)\propto a_i(t)^{-3}$ and
$\rho_{\mathrm{eff}}(t)=\frac{1}{N}\sum_i\rho_i(t)$, then generically
$\rho_{\mathrm{eff}}(t)\not\propto A(t)^{-3}$.
\end{lemma}

\begin{proof}
Let $y_i(t)=-3\ln a_i(t)$ so that $a_i(t)^{-3}=e^{y_i(t)}$ and
\[
A(t)^{-3}=\exp\!\left(\frac{1}{N}\sum_{i=0}^{N-1} y_i(t)\right).
\]
Since $\exp$ is strictly convex, Jensen's inequality gives
\[
\frac{1}{N}\sum_{i=0}^{N-1} e^{y_i(t)}
\;\ge\;
\exp\!\left(\frac{1}{N}\sum_{i=0}^{N-1} y_i(t)\right),
\]
with equality iff all $y_i(t)$ are equal, i.e.\ all $a_i(t)$ are equal.
\end{proof}

\section{Pedagogical applications}

\subsection{Classroom exercises}

\begin{exercise}[Sign mixing without global turning]
Construct $H_i(t)$ for $N=3$ such that $H_A(t)=0$ but $H_i(t)\neq 0$ for all $i$ at some $t$. Discuss what is physical and what is purely statistical.
\end{exercise}

\begin{exercise}[Closure sensitivity]
Compare three ``global scale'' definitions (geometric mean, arithmetic mean, and a volume-weighted mean). Show how inferred ``global curvature'' from an imposed closure depends on the definition, connecting to the fitting problem \citep{Buchert2000,ClarksonEtAl2011}.
\end{exercise}

\begin{exercise}[Dust mismatch computation]
Take $(a_0,a_1,a_2)=(1,2,4)$ at some fixed $t$. Compute $\langle a^{-3}\rangle$ and $A^{-3}$ and verify the lemma numerically.
\end{exercise}

\begin{exercise}[Effective-fluid completion]
Assume $\rho_X\propto a^{-n}$. Use \eqref{eq:wX} to derive $w_X=-1+n/3$. Explain why $\Lambda$-like behavior requires $n\approx 0$.
\end{exercise}

\begin{exercise}[Meta-structure brainstorming]
Propose a meta-structure that makes cross-universe time identification meaningful. For your proposal, list: (i) what selects the identification, and (ii) what degrees of freedom could mediate coupling.
\end{exercise}

\subsection{Research methodology implications}

\begin{itemize}[leftmargin=1.2em]
\item \textbf{Averaging requires a measure:} weights $w_i$ are not physics unless justified by a parent theory or observational principle \citep{Buchert2000,ClarksonEtAl2011}.
\item \textbf{Closure must be labeled:} postulates like \eqref{eq:globalClosure} should be stated explicitly as closure assumptions, not derived laws.
\item \textbf{Constraint modifications require completion:} Bianchi-consistency is non-negotiable; either complete via an effective fluid or derive from a covariant parent theory \citep{Wald1984}.
\end{itemize}

\section{Conclusion}

The ensemble construction shows how quickly cosmological narratives can be produced once global statistics are defined and a closure relation is postulated. The ``Eternal Bang'' story---simultaneous expansion and collapse within a cyclic meta-description---is best understood as a kinematic parable unless two gates are crossed: (i) a meta-structure that makes cross-universe time identification meaningful, and (ii) covariant (Bianchi-consistent) completion or derivation of any added coupling. These gates delineate the boundary between mathematical bookkeeping and physical cosmology.

\appendix

\section{Effective-fluid completion example}
Consider a single distinguished universe $U_0$ with scale factor $a(t)$ and Hubble rate $H=\dot a/a$. Suppose one writes a modified constraint
\[
H^2=\frac{8\pi G}{3}\left(\rho+\rho_X\right)-\frac{k c^2}{a^2}+\frac{\Lambda}{3}.
\]
Imposing $\nabla_\mu T_X^{\mu\nu}=0$ yields \eqref{eq:wX}. For $\rho_X=\alpha\,\rho_{\mathrm{ext}}$ one must specify $\rho_{\mathrm{ext}}(a)$ (or derive it from a parent theory) to determine the implied $w_X(a)$.

\section{Time-symmetric mirror pairs}
Assume $N$ even and pair $i\leftrightarrow\tilde i$ with
\[
a_{\tilde i}(t)=a_i(-t),\quad \rho_{\tilde i}(t)=\rho_i(-t),\quad p_{\tilde i}(t)=p_i(-t).
\]
Then
\[
H_{\tilde i}(t)=\frac{\dot a_{\tilde i}(t)}{a_{\tilde i}(t)}=\frac{-\dot a_i(-t)}{a_i(-t)}=-H_i(-t).
\]
If $H_i(0)=0$ and $\dot H_i(0)\neq0$, then for small $t>0$ one has $H_i(t)\approx \dot H_i(0)\,t$ and $H_{\tilde i}(t)\approx -\dot H_i(0)\,t$, giving explicit sign-mixing at the same $t$.

\section{Suggested assignment problems}

\begin{exercise}[Average-zero without individual zeros]
Show that $H_A(t)=0$ does not require any $H_i(t)=0$ by giving an explicit example for $N=3$.
\end{exercise}

\begin{exercise}[When does dust scaling survive?]
Assume $a_i(t)=c_i\,a(t)$ with constants $c_i>0$ independent of $t$. Show that
$\langle a^{-3}\rangle = C\,A^{-3}$ for an explicit constant $C$ depending only on $\{c_i\}$, and compute $C$.
\end{exercise}

\begin{exercise}[Consistent acceleration equation]
Starting from a modified constraint with $\rho_X$ and conservation of $\rho_X$, derive the corresponding acceleration equation and verify mutual consistency using the Bianchi identity logic \citep{Wald1984}.
\end{exercise}

\begin{exercise}[Explicit $A(t)$ in a mirror-pair toy example]
Take $N=2$ with $a_0(t)=e^{t^2}$ and $a_1(t)=a_0(-t)$. Compute $A(t)$ and $H_A(t)$ explicitly and interpret.
\end{exercise}

\begin{exercise}[Meta-structure proposals]
List three possible meta-structures that could justify a shared parameter $t$ across members, and for each describe what could mediate coupling.
\end{exercise}

\bibliographystyle{unsrt}
\bibliography{references}
\clearpage

\end{document}
